\section{Method}
\label{sec:method}

We first formalize a cross-resolution failure taxonomy derived from the
diagnosis (Section~\ref{sec:taxonomy}), then describe the Visibility-Boundary-Aware
Detection (VBAD) framework (3.2), its visibility head (3.3), and an ambiguity
boundary analysis (3.4).

\subsection{Cross-resolution failure taxonomy}
\label{sec:taxonomy}
Let a detector at input resolution $r$ produce a set of detections $D_r$ for an
image. A ground-truth defect $g$ is \emph{matched} by $D_r$ when some
detection $d \in D_r$ satisfies
\begin{equation}
\mathrm{match}(d, g) \;=\; \big[\,\mathrm{IoU}(d, g) \ge 0.5\,\big] \wedge
\big[\,\mathrm{cls}(d) = \mathrm{cls}(g)\,\big].
\end{equation}
Write $m_r(g) = 1$ if any detection in $D_r$ matches $g$, else $0$. Using a
low-resolution detector ($r{=}640$) and a high-resolution detector
($r{=}1280$), we partition every ground-truth defect into three types:
\begin{align}
\text{Type-N (normal)}:&\quad m_{640}(g) = 1,\\
\text{Type-V (visibility-limited)}:&\quad m_{640}(g) = 0 \;\wedge\; m_{1280}(g) = 1,\\
\text{Type-A (ambiguous)}:&\quad m_{640}(g) = 0 \;\wedge\; m_{1280}(g) = 0.
\end{align}
Type-N defects are already recovered cheaply. Type-V defects are the actionable
set: their visual evidence is insufficient at $640$ but sufficient at $1280$,
so they can be recovered by paying a higher resolution cost. Type-A defects
remain missed even at high resolution; as the per-class analysis in
Section~\ref{sec:exp} shows, these concentrate in extremely small
($<0.05\%$ image area) and weak-texture categories whose annotations are
themselves scale-sensitive, so we treat them as a reliability boundary rather
than an optimization target. Crucially, the Type-V label is obtained
automatically from cross-resolution failure agreement, not from manual
heuristics.

\subsection{VBAD framework}
\label{sec:vbad}
The key engineering question is \emph{when} to pay the high-resolution cost.
Running $1280$ on every image is expensive and, as we show in
Section~\ref{sec:exp}, can even degrade easy images due to scale-distribution
shift. VBAD therefore performs a per-image triage (Fig.~
\ref{fig:framework}):
\begin{enumerate}[leftmargin=1.4em]
\item Run the low-resolution detector to obtain detections $D_{640}$ and a
backbone feature descriptor $\phi$.
\item A lightweight \emph{visibility head} maps $\phi$ to a score
$s_{\mathrm{vis}} \in [0, 1]$ estimating whether the image contains Type-V
defects.
\item If $s_{\mathrm{vis}} \ge \tau$, run the high-resolution detector on the
\emph{full image} and use $D_{1280}$; otherwise keep $D_{640}$.
\end{enumerate}
We deliberately use full-image high-resolution reinspection rather than cropped
patches. A patch-based variant (Section~\ref{sec:exp}) recovers far fewer
defects because cropping discards the global context and scale distribution the
high-resolution detector relies on. Because Type-N and Type-A defects do not
benefit from the switch, and the two classes the framework routes between
(640 vs 1280 outputs) are taken as-is, the merge step is a simple selection and
needs no cross-resolution box fusion.

\subsection{Visibility head}
\label{sec:vishead}
The visibility head is trained as an image-level binary classifier. Its target
is
\begin{equation}
y_{\mathrm{vis}} = \mathbb{1}\big[\,\exists\, g : \text{Type-V}(g) = 1\,\big],
\end{equation}
i.e.\ the image contains at least one Type-V defect. The descriptor $\phi$
concatenates the globally pooled backbone feature of the $640$ detector with
four cheap scalar cues available from the same forward pass: the number of
low-confidence detections (score in $[0.05, 0.25)$, exposing
threshold-suppressed boxes), the number of confident detections, the smallest
predicted box area, and the mean detection confidence. A two-layer perceptron
maps $\phi$ to $s_{\mathrm{vis}}$, optimized with a class-balanced binary
cross-entropy loss
\begin{equation}
\mathcal{L} = \mathrm{BCE}\big(s_{\mathrm{vis}}, y_{\mathrm{vis}}; w_{+}\big),
\end{equation}
where $w_{+}$ reweights the positive class by its inverse frequency. The
detector is \emph{frozen} while the visibility head is trained, isolating the
triage-learning problem from detector optimization and keeping the added cost
of the head negligible relative to a detector forward pass.

\subsection{Ambiguity boundary analysis}
\label{sec:ambiguity}
Beyond triage, VBAD exposes a reliability boundary. The Type-A set---defects
missed even at high resolution---quantifies how many misses are
\emph{inherently unrecoverable} by resolution alone. We report the Type-A
fraction overall and per class, together with object-size and annotation
statistics, as a diagnostic rather than a trained component. In an industrial
inspection setting this boundary is directly actionable: images or classes with
high Type-A rates are precisely those an automated system should escalate to
manual review instead of optimizing against.
